%===========================================
% ESTE ARQUIVO É UMA INSTRUÇÃO/PROMPT DETALHADO PARA INTELIGÊNCIA ARTIFICIAL.
%
% OBJETIVO: Criar um template LaTeX COMPILÁVEL, TOTALMENTE MODULAR e livre de "gambiarras",
% utilizando EXCLUSIVAMENTE a classe MEMOIR para atender aos padrões ABNT e aos
% requisitos de numeração e estilo do usuário.
%
% OBSERVAÇÃO CRÍTICA: A modularização (criação de macros, \makechapterstyle, etc.)
% DEVE ser feita diretamente neste arquivo .tex, NÃO em um arquivo de classe (.cls)
% separado, para facilitar os testes e a depuração.
%===========================================
\documentclass[12pt,a4paper,oneside]{memoir}
\usepackage[brazilian]{babel}
\usepackage{fontspec}
\usepackage{tikz} % NECESSÁRIO PARA O ESTILO DE CAPÍTULO PERSONALIZADO
%===========================================

\setmainfont{CMU Serif}
\setlrmarginsandblock{3cm}{2cm}{*}
\setulmarginsandblock{3cm}{2cm}{*}
\checkandfixthelayout
\OnehalfSpacing % Espaçamento entre linhas de 1,5
\setlength{\parindent}{1.25cm} % Recuo do paragrafo
%===========================================

% ===========================================
% INSTRUÇÃO IA: CONFIGURAÇÃO DE ESTILO DE CAPÍTULO MODULAR (O SWITCH)
% ===========================================

% 1. Estilo Numerado Principal (Com TikZ) - Já fornecido
\makechapterstyle{CapitulosNumerado}{
    \setlength{\beforechapskip}{0.46cm}
    \renewcommand{\printchaptername}{}
    \renewcommand{\chapternamenum}{}
    \renewcommand{\afterchapternum}{\par\nobreak\vskip 1.3cm}
    \renewcommand{\printchapternum}{%
        \noindent
        \begin{tikzpicture}[remember picture, overlay]
            \def\numsize{61}
            \def\barwidth{0.8pt}
            \def\leftoffset{2.3cm}

            \node[inner sep=-1pt, anchor=south west, opacity=0] (num) at (\leftoffset, 0)
             {\fontsize{\numsize}{\numsize}\selectfont\rmfamily\thechapter};
            \node[inner sep=0pt, anchor=south west] at (\leftoffset, -1pt)
             {\fontsize{\numsize}{\numsize}\selectfont\rmfamily\thechapter};
            \node[inner sep=0pt, anchor=south west, font=\Large\rmfamily] (chaptext)
                 at (0pt, -4pt) {\chaptername};

            \coordinate (StartPoint) at (0.5*\barwidth, 13pt);
            \draw[line width=\barwidth] (StartPoint) -- (StartPoint |- num.north) -- (num.north west);
            \coordinate (RightWall) at (\linewidth, 0);
            \draw[line width=\barwidth] (num.north east) -- (RightWall |- num.north) -- (RightWall |- num.south) -- (num.south east);
        \end{tikzpicture}%
    }
    \renewcommand{\printchaptertitle}[1]{%
        \raggedright\huge\bfseries ##1
    }
}

% 2. Estilo Não Numerado Simplificado (ABNT: Centralizado com Reta) - CRIADO PELA IA
\makechapterstyle{estiloNaoNumerado}{
    \renewcommand{\printchaptername}{} % Sem nome do capítulo
    \renewcommand{\chapternamenum}{} % Sem número do capítulo
    \renewcommand{\afterchapternum}{\par\nobreak\vskip 1cm} % Espaço após o título
    \renewcommand{\printchapternum}{} % Sem número decorativo
    \renewcommand{\printchaptertitle}[1]{%
        \centering
        \vspace*{0.5cm} % Espaçamento superior
        \hrule
        \vspace*{0.5cm}
        \Huge\bfseries\MakeUppercase{##1} % Título em CAIXA ALTA
        \vspace*{0.5cm}
    }
}

% 3. Comando Modular (\capituloCentralizado) - CRIADO PELA IA
% Alterna para o estilo simplificado, aplica o capítulo* e reverte.
\newcommand{\capituloCentralizado}[2]{
    \chapterstyle{estiloNaoNumerado} % Ativa o estilo simplificado
    \chapter*{#1}
    \addcontentsline{toc}{chapter}{#2} % Adiciona ao sumário (Título em caixa alta)
    \chapterstyle{CapitulosNumerado} % Retorna ao estilo principal
}

%===========================================
% DEFINIÇÕES DE PÁGINAS DE ABERTURA
%===========================================
\newcommand{\folha}{
    \thispagestyle{empty}
    \vspace*{8cm}
    \begin{center}
        {\bfseries\scshape\Huge Folha de Rosto \par}
    \end{center}
    \clearpage
}
%===========================================
\newcommand{\contra}{
    \thispagestyle{empty}
    \vspace*{8cm}
    \begin{center}
        {\bfseries\scshape\Huge Contra Capa \par}
    \end{center}
    \clearpage
}
%===========================================

%===========================================
% INSTRUÇÃO IA: CONFIGURAÇÃO DE CABEÇALHOS (RUNNING HEADS) - IMPLEMENTADO PELA IA
%===========================================
% 1. Estilo de Página (\makepagestyle{meuestilo}) - Fornecido
\makepagestyle{meuestilo}
\makeevenhead{meuestilo}{\footnotesize\textit{\thechapter.}\quad\textit{\leftmark}}{}{\footnotesize\thepage}
\makeoddhead {meuestilo}{\footnotesize\textit{\thechapter.}\quad\textit{\leftmark}}{}{\footnotesize\thepage}
\makeheadrule{meuestilo}{\textwidth}{0.3pt}

% 2. Geração da Marca (\leftmark): Redefinido para incluir APÊNDICE/ANEXO/Número
% A classe memoir cuida da marcação \leftmark automaticamente para \chapter.
% O \thechapter será redefinido no \appendix/\annex.

% 3. Páginas de Abertura: A classe memoir usa \thispagestyle{plain} automaticamente na primeira página de \chapter.
%===========================================

%===========================================
% COMANDOS PÓS-TEXTUAIS CUSTOMIZADOS
%===========================================
% Redefine para começar o Apêndice em A e muda a nomenclatura
\newcommand{\meuAppendix}{
    \clearpage
    \appendix
    \renewcommand{\thechapter}{\Alph{chapter}} % Numeração em A, B, C...
    \renewcommand{\chaptername}{APÊNDICE} % Nome na TOC e no cabeçalho
    % Ajuste da marcação para o cabeçalho (ex: APÊNDICE A)
    \renewcommand{\memoirchaptermark}[1]{\markboth{\chaptername\ \thechapter.\ #1}{}}
}

% Redefine para começar o Anexo em A e muda a nomenclatura
\newcommand{\meuAnexo}{
    \clearpage
    \setchapterstyle{plain} % Reseta o estilo para evitar bugs de contador
    \setcounter{chapter}{0} % Reinicia o contador para começar de novo
    \renewcommand{\chaptername}{ANEXO} % Nome na TOC e no cabeçalho
    % Ajuste da marcação para o cabeçalho (ex: ANEXO A)
    \renewcommand{\memoirchaptermark}[1]{\markboth{\chaptername\ \thechapter.\ #1}{}}
}
%===========================================


\begin{document}

% ===========================================
% INÍCIO
% ===========================================

% 1. Numeração ARÁBICA (1, 2, 3...) (Invisível/Oculta na \frontmatter)
\pagenumbering{arabic}
\frontmatter % Ativa a numeração arábica, mas a classe memoir a oculta nas páginas \chapter*

\folha
\contra

%===========================================
% INSTRUÇÃO IA: ELEMENTOS PRÉ-TEXTUAIS
%===========================================
\capituloCentralizado{Agradecimento}{AGRADECIMENTO}
\capituloCentralizado{Resumo}{RESUMO}
\capituloCentralizado{Lista de Figuras}{LISTA DE FIGURAS}
\capituloCentralizado{Lista de Tabelas}{LISTA DE TABELAS}
\capituloCentralizado{Lista de Abreviaturas e Siglas}{LISTA DE ABREVIATURAS E SIGLAS}
\capituloCentralizado{Lista de Símbolos}{LISTA DE SÍMBOLOS}

%===========================================
% INSTRUÇÃO IA: CONFIGURAÇÃO DO SUMÁRIO
%===========================================
\renewcommand{\cfttoctitlefont}{\hfill\Huge\bfseries\MakeUppercase} % Centraliza e estiliza o título
\renewcommand{\cftaftertoctitle}{\hfill}
\setlength{\cftchapnumwidth}{2cm} % 2 cm de espaçamento horizontal para números
\settocstyle{allcaps} % Títulos em CAIXA ALTA
\tableofcontents* % Supressão do Título "SUMÁRIO"
%===========================================


%===========================================
% INSTRUÇÃO IA: TRANSIÇÃO PARA A PARTE TEXTUAL (CRÍTICO)
%===========================================
\mainmatter % 1. Ativa a exibição VISÍVEL da numeração arábica

\pagestyle{meuestilo} % 2. Ativa o estilo de cabeçalho 'meuestilo'

\chapterstyle{CapitulosNumerado} % 3. Ativa o estilo de capítulo com TikZ (Para Capítulos Numerados)

% Introdução (Híbrida: \chapter* mas na \mainmatter)
\capituloCentralizado{Introdução}{INTRODUÇÃO}

% Contagem de Capítulos DEVE continuar sequencialmente (1, 2, 3...).
\part{PREPARAÇÃO DO RELATÓRIO}
\chapter{Desenvolvimento do Tema}
\part{RESULTADO DO RELATÓRIO}
\chapter{Conclusão do Trabalho}


%===========================================
% INSTRUÇÃO IA: ELEMENTOS PÓS-TEXTUAIS (CRÍTICO)
%===========================================
\backmatter % 1. Transição Pós-Textual (\chapter* não é numerado no \backmatter, e a numeração continua)

% 2. Referências: \chapter*
\capituloCentralizado{Referência}{REFERÊNCIA}

% Apêndices
\meuAppendix % 4. Numeração: Reseta o contador, define o estilo e a nomenclatura.
% A IA deve criar um comando que use \chapter* + \addcontentsline{toc} sem número de página
\capituloCentralizado{Apêndices}{\protect\addtocontents{toc}{\protect\vskip 1em}\protect\nopagebreak[4]\protect\centering APÊNDICES\protect\par} % Título divisório no TOC

\chapter{Título do Apêndice A} % Deve gerar APÊNDICE A
\chapter{Título do Apêndice B} % Deve gerar APÊNDICE B

% Anexos
\meuAnexo % 4. Numeração: Reseta o contador, define o estilo e a nomenclatura.
% A IA deve criar um comando que use \chapter* + \addcontentsline{toc} sem número de página
\capituloCentralizado{Anexos}{\protect\addtocontents{toc}{\protect\vskip 1em}\protect\nopagebreak[4]\protect\centering ANEXOS\protect\par} % Título divisório no TOC

\chapter{Título do Anexo C} % Deve gerar ANEXO A (contador reiniciado)
\chapter{Título do Anexo D} % Deve gerar ANEXO B (contador reiniciado)

\end{document}